% !TEX root = ../../ctfp-print.tex

\lettrine[lhang=0.17]{ま}{だ圏論がすべて射についてのものだと}納得させられていないなら、私は自分の役目を十分に果たせていないことになる。次のトピックは随伴（adjunctions）で、これはhom集合の同型によって定義される。そのため、hom集合の構成要素についての直感を見直しておくことには意味がある。また、随伴はこれまで学んできた多くの構成を記述するより一般的な言語を提供してくれるので、それらを復習しておくことも役に立つかもしれない。

\section{関手}

まず、関手は射の写像として考えるべきだ。これはHaskellにおける\code{Functor}型クラスの定義で強調されている見方であり、\code{fmap}を中心に回っている。もちろん、関手は対象も写す。射の端点を写さなければ、合成を保存することについて語れないからだ。対象は、どの射の対が合成可能かを教えてくれる。ある射のターゲットが別の射のソースと等しくなければ、それらは合成できない。したがって、射の合成が\newterm{持ち上げられた}（lifted）射の合成に写されるようにしたいなら、それらの端点の写像はほぼ自動的に決まる。

\section{可換図式}

射の多くの性質は、可換図式によって表現される。ある射が他の射の合成として複数の方法で記述できる場合、可換図式が成立する。

特に、可換図式はほぼすべての普遍的構成の基礎をなす（始対象と終対象は注目すべき例外だ）。積、余積、その他さまざまな（余）極限、指数対象、自由モノイドなどの定義でこれを見てきた。

積は普遍的構成の簡単な例だ。二つの対象$a$と$b$を選び、射の対$p$と$q$とともに対象$c$が存在して、それらの積の普遍性を持つかどうかを確認する。

\begin{figure}[H]
  \centering
  \includegraphics[width=0.3\textwidth]{images/productranking.jpg}
\end{figure}

\noindent
積は極限の特殊ケースだ。極限は錐によって定義される。一般的な錐は可換図式から構築される。それらの図式の可換性は、関手の写像に対する適切な自然性条件で置き換えることができる。こうして可換性は、自然変換というより高レベルの言語のためのアセンブリ言語の役割に還元される。

\section{自然変換}

一般に、自然変換は射から可換正方形への写像が必要なときに非常に便利だ。自然性の正方形の向かい合う二辺は、二つの関手$F$と$G$のもとでの射$f$の写像だ。残りの二辺は自然変換の成分（これらも射だ）である。

\begin{figure}[H]
  \centering
  \includegraphics[width=0.35\textwidth]{images/3_naturality.jpg}
\end{figure}

\noindent
自然性とは、「隣接する」成分に移動するとき（隣接とは射によって結ばれていることを意味する）、圏または関手のどちらの構造にも逆らわないことを意味する。対象間のギャップを埋めるために自然変換の成分を使ってから関手で隣に飛ぶか、あるいはその逆の順序で行うかは関係ない。二つの方向は直交している。自然変換は左右に移動させ、関手は上下あるいは前後に移動させる、とでも言えばよいか。関手の\emph{像}（image）をターゲット圏のシートとして可視化できる。自然変換は$F$に対応するシートを、$G$に対応する別のシートに写す。

\begin{figure}[H]
  \centering
  \includegraphics[width=0.35\textwidth]{images/sheets.png}
\end{figure}

\noindent
Haskellでこの直交性の例を見てきた。そこでは関手の作用はコンテナの形を変えずにその内容を変更し、一方で自然変換は触れていない内容を別のコンテナに詰め替える。これらの操作の順序は関係ない。

極限の定義における錐が自然変換で置き換えられることも見てきた。自然性はすべての錐の辺が可換であることを保証する。それでも、極限は錐\emph{間}の写像によって定義される。これらの写像も可換性条件を満たさなければならない（たとえば、積の定義における三角形は可換でなければならない）。

これらの条件も自然性によって置き換えることができる。\emph{普遍的な}錐、つまり極限は、（反変）hom関手：
\[F \Colon c \to \cat{C}(c, \Lim[D])\]
と、対象を$\cat{C}$の錐（それ自体が自然変換だ）に写す（同じく反変の）関手：
\[G \Colon c \to \mathit{Nat}(\Delta_c, D)\]
の間の自然変換として定義されることを思い出してほしい。ここで、$\Delta_c$は定数関手、$D$は$\cat{C}$における図式を定義する関手だ。関手$F$と$G$はともに$\cat{C}$の射への well-defined な作用を持つ。この$F$と$G$の間の特定の自然変換が\emph{同型}であることがわかる。

\section{自然同型}

\newterm{自然同型}（natural isomorphism）とは、すべての成分が可逆な自然変換のことだ。これは「二つのものが同じ」という圏論の言い方だ。そのような変換の成分は対象間の同型、つまり逆を持つ射でなければならない。関手の像をシートとして可視化するなら、自然同型はそれらのシート間の一対一の可逆写像だ。

\section{Hom集合}

では、射とは何だろうか？射は対象よりも多くの構造を持つ。対象と違い、射は二つの端点を持つ。しかし、ソース対象とターゲット対象を固定すると、二者間の射は退屈な集合を形成する（少なくとも局所的に小さい圏では）。この集合の要素に$f$や$g$といった名前を付けて区別することができるが、実のところ何がそれらを異ならせているのだろうか？

与えられたhom集合における射の本質的な違いは、隣接するhom集合からの他の射とどのように合成されるかにある。射$h$があって、$f$との合成（前合成でも後合成でも）が$g$との合成と異なる場合、たとえば：
\[h \circ f \neq h \circ g\]
ならば、$f$と$g$の違いを直接「観測」できる。しかし違いが直接観測できない場合でも、関手を使ってhom集合を拡大することができる。関手$F$は二つの射を異なる射に写すかもしれない：
\[F f \neq F g\]
これはより豊かな圏の中でのことで、隣接するhom集合がより高い解像度を提供する。たとえば：
\[h' \circ F f \neq h' \circ F g\]
ここで$h'$は$F$の像の中にない。

\section{Hom集合の同型}

多くの圏論的構成は、hom集合間の同型に依存している。しかしhom集合は単なる集合なので、それらの間の単純な同型はあまり多くを語らない。有限集合の場合、同型は単に同じ数の要素を持つことを意味する。集合が無限の場合、その濃度が等しくなければならない。しかし、hom集合の意味のある同型は合成を考慮しなければならない。合成は複数のhom集合を伴う。hom集合のコレクション全体にわたる同型を定義し、合成と相互作用する互換性条件を課す必要がある。そして\newterm{自然}（natural）同型がまさにそれを満たす。

では、hom集合の自然同型とは何だろうか？自然性は関手間の写像の性質であり、集合間の写像の性質ではない。つまり、hom集合値関手間の自然同型について話しているのだ。これらの関手は単なる集合値関手以上のものだ。射への作用は適切なhom関手によって誘導される。射はhom関手によって前合成または後合成（関手の共変性に応じて）を用いて標準的に写される。

米田埋め込みはそのような同型の一例だ。$\cat{C}$のhom集合を関手圏のhom集合に写す。しかも自然に。米田埋め込みにおける一方の関手は$\cat{C}$のhom関手で、もう一方は対象をhom集合間の自然変換の集合に写す。

極限の定義もhom集合間の自然同型だ（二番目のものは再び関手圏にある）：
\[\cat{C}(c, \Lim[D]) \simeq \mathit{Nat}(\Delta_c, D)\]
指数対象の構成や自由モノイドの構成も、hom集合間の自然同型として書き直せることがわかる。

これは偶然ではない。次に見るように、これらはすべて随伴の異なる例にすぎず、随伴はhom集合の自然同型として定義される。

\section{Hom集合の非対称性}

随伴を理解するのに役立つもう一つの観察がある。Hom集合は一般に対称ではない。hom集合$\cat{C}(a, b)$はhom集合$\cat{C}(b, a)$と大きく異なることが多い。この非対称性の究極の実例は、圏として見た半順序集合だ。半順序では、$a$から$b$への射は$a$が$b$以下の場合にのみ存在する。$a$と$b$が異なるならば、逆方向の射、つまり$b$から$a$への射は存在できない。したがって、hom集合$\cat{C}(a, b)$が空でない（この場合は単集合を意味する）ならば、$a = b$でない限り$\cat{C}(b, a)$は空でなければならない。この圏の矢印は一方向に定まった流れを持つ。

前順序は、必ずしも反対称でない関係に基づく順序で、時折サイクルが存在することを除いて「ほとんど」方向性を持つ。任意の圏を前順序の一般化として考えると便利だ。

前順序は薄い圏（thin category）であり、すべてのhom集合は単集合か空集合のどちらかだ。一般的な圏は「厚い」前順序として可視化できる。

\section{練習問題}

\begin{enumerate}
  \tightlist
  \item
        自然性条件のいくつかの退化したケースを考え、適切な図式を描け。たとえば、関手$F$または$G$のどちらかが対象$a$と$b$（$f \Colon a \to b$の端点）の両方を同じ対象に写す場合、たとえば$F a = F b$または$G a = G b$ならば何が起きるか？（このようにして錐または余錐が得られることに注目せよ。）次に$F a = G a$または$F b = G b$の場合を考えよ。最後に、自分自身へのループ射$f \Colon a \to a$から始めたらどうなるか？
\end{enumerate}
